%!TEX root = ./../../springer-mono-format/book.tex
%!TEX root = ./../../thesis-trier.tex

\IN{binary relation} \label{def:relation}
\IN{relation} 
A \emph{relation} $R$ is a tuple $(R_G, X_1, \ldots, X_n$)
such that $R_G \subseteq X_1 \times \cdots \times X_n$.
For $n = 2$, we call $R$ a \emph{binary relation}
and in this case, for $(x_1, x_2) \in X_1 \times X_2$
we write $x_1 R x_2$ if and only if $(x_1, x_2) \in R_G$.

\IN{homogeneous relation}
For $X = X_1 = X_2$ we say that $R$ is a relation over $X$ (such relations are also called \emph{homogeneous}). 
If~$R$ is a relation over $X$,
then it is called
\begin{description}  
	
	\IN{reflexive relation}
	\item[\normalfont{\emph{reflexive}}] if for all $x \in X$ we have $xRx$,
	
	\IN{symmetric relation}
	\item[\normalfont{\emph{symmetric}}] if for all $x,y \in X$ we have that $xRy$ implies $yRx$, 
	
	\IN{antisymmetric relation}
	\item[\normalfont{\emph{antisymmetric}}] if for all $x,y \in X$ we have that $xRy$ and $yRx$ imply $x = y$, 
	
	\IN{transitive relation}
	\item[\normalfont{\emph{transitive}}] if for all $x,y,z \in X$ we have that $xRy$ and $yRz$ imply $xRz$,
	
	\IN{equivalence relation}
	\item[\normalfont{\emph{equivalence relation}}] if it is reflexive, symmetric and transitive. 
\end{description}

\label{def:eqfactor}
\IN{equivalence class}
If $\sim$ denotes an equivalence relation\footnote{Using symbols like $\sim$ for equivalence relations is more common than denoting them by~$R$.}
on $X$ and $x \in X$,
the \emph{equivalence class of~$x$} is 
\[
[x]_{\sim} = \{ y \in X \mid y \sim x \}. \label{def:equivalenceclasses}
\]
It is $[x]_{\sim} = [y]_{\sim}$ if and only if $x \sim y$ for $x,y \in X$.
We write $X / {\sim} = \{ [x]_{\sim} \mid x \in X \}$ 
for the set of equivalence classes. The surjective map $\pi_{\sim} \colon X \to X / {\sim}$ given by
\begin{equation} \label{def:canprojmap}
 \pi_{\sim}(x) = [x]_{\sim}
\end{equation}
for $x \in X$ is called the \emph{canonical projection map (for $\sim$)}. 
\IN{canonical projection map} 

\IN{equivalence relation induced!by partition}
The equivalence classes partition the set $X$. Conversely, given a partition $\{ X_i \mid i \in I\}$ of~$X$
for some index set $I$, the \emph{equivalence relation induced by it} is given by setting $x,y \in X$ to be equivalent
if $\{x,y\} \subseteq X_i$ for some $i \in I$. 

For a function $f \colon X \to Y$
the non-empty preimages $\applyfunc{y}{f^{-1}}$, $y \in Y$,
partition $X$. The equivalence relation $\sim_f$ \label{def:kernel}
induced by this partition is called the \emph{kernel relation~(of~$f$)}. We have 
$x \sim_f y$ if and only if $\applyfunc{x}{f} = \applyfunc{y}{f}$.
The kernel relation of the canonical projection $\pi_{\sim}$ from above is $\sim$.
\IN{kernel relation}

\IN{saturating a subset}
We say that an equivalence relation $\sim$ on $X$ \emph{saturates} 
a subset $A \subseteq X$ if and only
if $x \sim y$ and $x \in A$ imply $y \in A$, or equivalently $A$ is union of equivalence classes.

\begin{proposition} \label{prelim:prop:equivalence-saturates}
	Let $f \colon X \to Y$ and $A \subseteq X$. 
	Then the following are equivalent:
\begin{enumerate}
\item \label{equivalence-saturates:item:simf} The kernel relation $\sim_f$ saturates $A$,

\item \label{equivalence-saturates:item:subset} $\applyfunc{\applyfunc{A}{f}}{f^{-1}} \subseteq A$,

\item \label{equivalence-saturates:item:eq} $\applyfunc{\applyfunc{A}{f}}{f^{-1}} = A$.
	
\end{enumerate}
\end{proposition}
\begin{proof}	
$\ref{equivalence-saturates:item:simf} \Rightarrow~\ref{equivalence-saturates:item:subset}$: That $\sim_f$ saturates $A$
 is equivalent to $f^{-1}(f(x)) \subseteq A$ for each $x \in A$.
 So $f^{-1}(f(A)) = f^{-1}(f(\bigcup_{x \in A} \{x\})) \subseteq A$.
 

$\ref{equivalence-saturates:item:subset} \Rightarrow~\ref{equivalence-saturates:item:eq}$: This 
follows from \cref{prelim:eqn:finvfinclusion}.
 
 
$\ref{equivalence-saturates:item:eq} \Rightarrow~\ref{equivalence-saturates:item:simf}$: If $x \sim_f y$ for some $x \in A$ and $y \in X$,
then $y \in f^{-1}(f(x)) \subseteq f^{-1}(f(A)) = A$.	
\end{proof} 

\IN{trivial equivalence relations}
\IN{equality relation}
\IN{universal relation}
The \emph{trivial equivalence relations} 
are $X \times X$ (the \emph{universal relation}) 
and $\{ (x,x) \mid x \in X\}$ (the \emph{equality relation}).

\IN{finer equivalence relation}
\IN{coarser equivalence relation}
\IN{generated!equivalence relation}
For two equivalence relations $\sim$ and $\approx$ 
on $X$ such that $x \sim y$ implies $x \approx y$ for all $x,y$, the relation~$\sim$ (resp.~$\approx$)
is said to be \emph{finer} (resp.~\emph{coarser}) 
than $\approx$ (resp.~$\sim$).
The \emph{equivalence relation generated}
by a given relation $R$ over $X$
is the finest equivalence relation $\sim$
such that $xRy$ implies $x \sim y$ for all $x,y \in X$, i.e.,
the unique relation such that
\begin{enumerate}
	\item $xRy$ implies $x \sim y$ for all $x,y \in X$ and
	
	\item for each equivalence relation $\equiv$ over $X$ such that 
	$xRy$ implies $x \equiv y$ for all $x,y \in X$ we have
	that $x \sim y$ implies $x \equiv y$ for all $x,y \in X$.
\end{enumerate}

\begin{remark} \label{prelim:rel:rem:def}
	In the literature, a relation $R$ from $X$ to $Y$ is often just defined
	as a subset of $X \times Y$. However, as pointed out in~\cite[p.~2]{Eilenberg1974},
	under the usual identification of $(X \times Y) \times Z$
	and $X \times (Y \times Z)$,
	a subset of $X \times Y \times Z$
	can be interpreted as a relation 
	from $X \times Y$ to $Z$
	or from $X$ to $Y \times Z$,
	and both interpretations may yield
	relations with quite different properties.	
	Another reason is that there might be (type- or category-theoretical) contexts
	when relations $(R, \emptyset, X)$
	and $(R', \emptyset, Y)$ for $X \ne Y$
	should be treated differently.
	Note that for homogeneous relations,
	these issues do not arise and the relation can indeed
	be identified with its set of tuples.
\end{remark}

\begin{remark}
	In~\cite{Eilenberg1974}	
	a relation from $X$ to $Y$ 
	is modeled as a function $f \colon X \to \mathcal P(Y)$,
	which can be identified with functions $f \colon \mathcal P(X) \to \mathcal P(Y)$
	with the property $f(\bigcup_{i\in I} S_i) = \bigcup_{i \in I} f(S_i)$
	for $S_i \in \mathcal P(X)$ with $i \in I$.	
	Note that, by the identities from \cref{prelim:subsec:func:basic},
	the extension of a function $X \to Y$ to subsets
	fulfills this equation.
\end{remark}